% =============================================================================
\section{Fix Verification with a Live Negative}
\label{sec:fix}

\subsection{What the fix actually changes}
\label{sec:fixchange}

The fix commit is public: \texttt{d7bc52be\allowbreak eadff4be5f5690de4d5de42a\allowbreak bd10affe},
dated 2026-07-26T17:32:02Z, titled ``refactor: git patch apply (\#38637)
(\#38638)'', touching \texttt{cherry\_pick.go}, \texttt{patch.go},
\texttt{patch\_test.go} and \texttt{update.go}. Inside that refactor sits one
decisive hunk (Listing~\ref{lst:fix}). The regression test added in the same
commit asserts only that \texttt{.git} exists inside the temporary repository:
the test encodes ``the patch sandbox must not be bare'' and nothing more.

\begin{lstlisting}[style=diffsrc,
caption={The decisive hunk of the fix (\texttt{services/repository/files/patch.go};
full diff in \texttt{designs/fix-commit-full.diff}). One boolean turns the patch
sandbox from bare to non-bare. The \texttt{git apply} flags of
Listing~\ref{lst:apply} are not touched.},
label={lst:fix}]
-  if err := t.Clone(ctx, opts.OldBranch, true); err != nil {
+  // here must NOT use bare repo, because the following git
+  // commands might operate working tree ("--index") directly
+  if err := t.Clone(ctx, opts.OldBranch, false); err != nil {
       return nil, err
  }
\end{lstlisting}

\textbf{The consequence that is easy to miss.} The fix changes \emph{where} a
patch is applied. It does not change the \texttt{git apply} flags; it adds no
path deny-list for \texttt{hooks/}, no filter on executable modes, and does not
disable the three-way fallback. The collision mechanism is still there --- the
sandbox simply no longer makes its checkout root double as
\texttt{\$GIT\_DIR}. Two operational consequences follow.

\begin{enumerate}
\item \textbf{A byte-identical exploit is inert, and still looks the same.} The
exploit bytes are not rejected, not sanitised, and not made anomalous; they are
applied to a worktree, where \texttt{hooks/post-index-change} is a plain file.
We confirmed this by observation rather than inference: on the patched instance
the monitor still saw the hook file materialise on disk
(Table~\ref{tab:ab}).
\item \textbf{Content-signature detection is therefore structurally weak.} A
signature over the patch text --- a \texttt{hooks/} path, \texttt{new file mode
100755}, the \texttt{diff} header --- fires identically on a vulnerable and a
patched instance, and fires on any legitimate patch that touches hook files. It
can raise an attempt signal; it cannot tell an operator whether they were
compromised, and it cannot tell them whether they are patched. What \emph{can}
discriminate is behaviour: process lineage, and which directory Git considered
\texttt{\$GIT\_DIR}. That is the most actionable claim in this paper, and it
comes from reading the fix rather than the exploit.
\end{enumerate}

\subsection{A/B experiment}

Same lab network, same configuration, same account shape, byte-identical payload;
the only variable is the image (\texttt{1.27.0} against \texttt{1.27.1}). The
negative arm is not accepted on absence of evidence: a polling monitor
(\texttt{exploit/ab-monitor.sh}) searched for the materialised hook file under
\texttt{/data/gitea/tmp/local-repo/} throughout each patched run, so that ``no
execution'' cannot be quietly explained by ``the payload never arrived''.

\begin{table}[t]
\centering
\footnotesize
\begin{tabular}{@{}p{2.7cm}p{1.9cm}p{2.9cm}@{}}
\toprule
Observation & \texttt{1.27.0} vuln. & \texttt{1.27.1} patched \\
\midrule
\texttt{diffpatch} 1/2 & HTTP 201 & HTTP 201 \\
\texttt{diffpatch} 2/2 & HTTP 201 & \textbf{HTTP 500} \\
Hook file at the sink path & yes & \textbf{observed} on disk, in a temporary
\texttt{upload.git*} repository \\
Canary created & yes & \textbf{no --- 22/22 cycles} \\
Hook content executed & yes, \texttt{uid=1000} & none observed \\
Delivery window & $0.82$--$0.89$\,s & not applicable \\
\bottomrule
\end{tabular}
\caption{A/B result. Identical payload and configuration. The patched arm is a
\emph{live} negative: the payload reaches the sink path and is still not
executed, and the second submission now fails with HTTP 500 instead of returning
a successful commit.}
\label{tab:ab}
\end{table}

\begin{figure}[t]
\centering
\includegraphics[width=\columnwidth]{\figdir/fig-still-vulnerable-arm.png}
\caption{The victim repository on the \emph{vulnerable} arm, mid-defacement on
the filmed run: the page's own footer reads \texttt{Version:~1.27.0} and the
address bar reads \texttt{172.31.224.2}; the crop keeps the full footer band,
its geometry measured in \texttt{paper/FIGURES.md}. The README has been
overwritten through the same primitive that landed the shell --- the table
visible here is the agent's own impact note, declaring which value was
replaced and how the original bytes are restored. The patched arm is not shown
on camera --- the demonstration film was ruled offensive-only --- so this
figure identifies only the arm it shows, and the patched-arm case rests on the
A/B table above and \texttt{evidence/fix-verification/}, not on a frame.}
\label{fig:arm}
\end{figure}

\textbf{Why the negative is non-vacuous.} Three observations have to hold
together for the patched arm to mean anything, and all three do. Fig.~\ref{fig:arm}
shows the arm under test as the operator saw it. (i) The payload
reaches the target and is written: the monitor captured
\texttt{hooks/post-index-change} under a temporary upload repository on the
\emph{patched} instance. (ii) The failure is visible in the protocol: the second
submission returns HTTP 500 where the vulnerable instance returns HTTP 201.
(iii) Nothing executes across 22 independent cycles --- we ran the chain 22 times
rather than once, because a single negative proves only that one run did not
fire. Given a non-bare sandbox, the checked-out hook lands in
\texttt{<worktree>/hooks/}, which is not \texttt{\$GIT\_DIR/hooks/}, so Git has no
reason to execute it. That is exactly the failure mode the hunk in
Listing~\ref{lst:fix} predicts, and that is the sense in which the fix is
verified: not ``it fails'', but ``it fails in the way, and only in the way, that
the change implies''.

We do not claim to have identified which internal call produces the HTTP 500; we
report the observed status code, not a mechanism we did not measure.

\subsection{Severity, checked against what we could prove}
\label{sec:cvss}

The severity of record is the NVD CVSS 3.1 assessment
\texttt{AV:N/AC:L/PR:N/UI:N/S:U/C:H/I:H/A:H} = 9.8 (Critical), typed
\emph{Secondary}, CWE-94. Our measurements are consistent with it: the endpoint
is network-reachable, no privilege beyond a self-created account is needed while
registration is open (Section~\ref{sec:precond-cred}), no user interaction is
involved, and total delivery took under a second --- so \texttt{AC:L} is not a
generous reading. Two honesty notes. First, \texttt{PR:N} rides entirely on open
registration: with registration disabled the prerequisite becomes one ordinary
authenticated account, effectively \texttt{PR:L}, and the vulnerability remains
exploitable --- the bar is ``any user who can create a repository'', not ``an
administrator''. Second, we did not test availability impact; our payload policy
forbids destructive action, so \texttt{A:H} is an inference from arbitrary write
access as the service user, not a result we produced.

% =============================================================================
\section{What a Defender Can Actually See}
\label{sec:detect}

This section states detection \emph{requirements} derived from measured
artifacts, and it quantifies the shipped rule package (\texttt{rules/}) against
captures we recorded ourselves. We publish no detection rate and no false-positive
rate that we did not measure: every figure below is re-derivable from the fire
reports in \texttt{evidence/rule-fires/}. What we can state is which signals
discriminate and which do not --- and that distinction is the useful part.
The sensor itself is evidenced by the fire reports and the tables below, not by
a screenshot: what we publish per signal is the command that ran, the artifact
it consumed, and the alert lines it produced.

\subsection{Measured behaviour of the shipped signatures}

\textbf{What counts.} The package carries seven CVE SIDs, of which \emph{six}
alert: \texttt{sid:202660045} is a \texttt{flowbits:\,set} carrying
\texttt{noalert}, so it can never appear in an alert count, and a harness that
reports ``seven signatures fired'' is wrong by construction. Replayed through
Suricata~8.0.6, the attack capture yields \textbf{11 alerts} from the six
alerting SIDs (\texttt{41:2}, \texttt{42:2}, \texttt{43:2}, \texttt{44:1},
\texttt{46:2}, \texttt{47:2}); both captures are \texttt{suricata -r} replays of
recorded wire data, not live sensor output.

\textbf{False positives, measured.} A benign negative is only worth what its
traffic is worth. Our simple capture (\texttt{wire-benign-control.pcap},
48\,packets) contains no request with the shape of the attack, so its zero
CVE-SID alerts prove that the engine decoded the bytes and nothing about
specificity --- reading it as a specificity result is the error we corrected
while preparing this paper. We therefore recorded a structurally similar control,
\texttt{wire-benign-structural.pcap} in that same folder (123 packets read),
whose traffic is what a legitimate administrator generates: three
\texttt{diffpatch} cycles that add an executable \texttt{maintenance/helper.sh},
built with the same patch generator the shipped PoC uses, plus ordinary
\texttt{git-upload-pack}\allowbreak/\allowbreak\texttt{git-receive-pack} traffic.

On that benign traffic the mode-bit signature \texttt{sid:202660041} fired on
\textbf{3 of 3} legitimate cycles, while the five signatures keyed to the
\texttt{hooks/} target path produced no alert. One of the six alerting
rules, on one capture, is not a measurement of false-positive rate in the field;
it is a measurement of how little a mode bit discriminates. That result is why
\texttt{sid:202660041} ships at \texttt{priority:3} with
\texttt{confidence~low} and its measured rate recorded in rule metadata, and why
we do not recommend it as a standalone alert.

\subsection{Signals that do not discriminate}

\textbf{The patch content.} Two \texttt{POST} requests to the
\texttt{diffpatch} route carrying identical bodies, the second declaring a
\texttt{new file mode 100755} under \texttt{hooks/}, are the shape of this
attack on the wire, and our capture (\texttt{evidence/wire-attack.pcap}, 33\,400
bytes, 72 packets, filtered on \texttt{tcp port 3000}) contains it alongside
benign control traffic --- a \texttt{GET /} and one ordinary \texttt{POST
contents} --- recorded precisely so that false-positive behaviour can be
measured. But per Section~\ref{sec:fixchange} the same bytes are accepted and
written on a patched instance, so a content signature is an \emph{attempt}
indicator with no bearing on outcome.

\textbf{The HTTP response.} Both \texttt{diffpatch} calls return HTTP 201 on a
vulnerable instance while code executes, because the hook's exit status is not
propagated. Success-code monitoring sees nothing. Note the asymmetry we
measured: it is the \emph{patched} instance that answers 500, so a ``many 5xx on
\texttt{diffpatch}'' alert is a statement about the fix, not about compromise.

\textbf{Disk forensics after the fact.} The hook lives in a short-lived
temporary upload repository, so the artifact that proves execution is gone by
the time an operator looks (Section~\ref{sec:impact}).

\subsection{Signals that do discriminate}

\begin{enumerate}
\item \textbf{Process lineage.} The parent of the hook shell is
\texttt{/usr/bin/git write-tree} and the child is \texttt{/bin/sh
hooks/post-index-change 0 0}, under the \texttt{gitea web} process, as
\texttt{uid=1000} (Listing~\ref{lst:canary}). A \texttt{git} command that writes
a tree should not have a shell descendant. This is the highest-fidelity signal we
measured: across the 394 service-log lines covering our run, not one mentions the
executing hook while 40 name the \texttt{diffpatch} endpoint itself, so the
request is visible and the execution is not
(\texttt{evidence/app-log-visibility.txt}); and it fires only when execution
actually happened.
\item \textbf{Location, not existence, of the hook file.} A file named
\texttt{hooks/post-index-change} inside a directory that \emph{is}
\texttt{\$GIT\_DIR} is compromise; the same file inside a worktree is inert. The
discriminating observable is therefore \texttt{GIT\_DIR=.} together with the
temporary path (\texttt{/data/gitea/tmp/local-repo/}, name prefix
\texttt{upload.git}), not the filename. A rule keyed on the filename alone
produces the same alert on patched and vulnerable hosts.
\item \textbf{Repeated identical \texttt{diffpatch} against one repository.} The
exploit \emph{requires} the same patch twice inside one second
(Table~\ref{tab:timing}). That repetition is a behavioural invariant of the
attack, independent of the file it carries.
\item \textbf{A new ref in a repository with no push.} This is the vendor's
retrieval channel: output committed as an object, surfaced as a branch, fetched
over authenticated smart HTTP. We did not reproduce it, so we list it as a
requirement the advisory implies rather than a signal we validated.
\end{enumerate}

\subsection{MITRE ATT\&CK mapping of what we demonstrated}

Table~\ref{tab:attack} is deliberately narrow: every row names a technique and the
concrete observation that backs it in this paper, with a section reference. Rows we
could not back are absent rather than inferred. In particular there is no Lateral
Movement row and no Exfiltration row, because we did not reproduce the vendor's
output-free retrieval channel and did not move laterally in the lab; we write a
\emph{requirement} the advisory implies in Section~\ref{sec:detect} instead of a
technique we observed. Mapping is at tactic level only: we make no claim about which
sub-technique an intrusion would be attributed to in the field.

\begin{table}[t]
\centering
\footnotesize
\begin{tabular}{@{}p{1.3cm}p{2.0cm}p{4.2cm}@{}}
\toprule
Tactic & Technique & Evidence in this paper \\
\midrule
Initial Access & T1190 & Self-registered non-admin account; HTTP 303 on
\texttt{sign\_up} (\S\ref{sec:precond-cred}) \\
Execution & T1059.004 & \texttt{/bin/sh hooks/}\allowbreak\texttt{post-index-change
0 0}, \texttt{uid=1000(git)} (Listing~\ref{lst:canary}) \\
Discovery & T1083 & \texttt{app.ini} and the repository tree enumerated from
inside the hook (\S\ref{sec:impact}) \\
Credential Access & T1552.001 & 3 secret-bearing lines in \texttt{app.ini};
planted secrets read back \\
Persistence & T1546.004 & Git hook installed through repository-controlled
content (lab only) \\
\bottomrule
\end{tabular}
\caption{Techniques mapped only to behaviour we actually observed. Lateral
Movement and Exfiltration are absent on purpose: the vendor's output-free
exfiltration channel was not reproduced here.}
\label{tab:attack}
\end{table}
